\section{Experiments}
\label{sec:experiments}

\subsection{Setup}
\label{sec:setup}

\noindent\textbf{Dataset.} We use
TotalSegmentator~\cite{wasserthalTotalSegmentatorRobustSegmentation2023}
CT, \TODO{state the train/held-out anatomical class split, following the
same held-in/held-out protocol used in our 2D work to test in-context
generalization to unseen classes, not just unseen subjects}.

\noindent\textbf{Baselines.} Medverse~\cite{huMedverseUniversalModel2025a}
and Iris~\cite{gaoShowSegmentUniversal2025}, \TODO{state whether
baselines are retrained on the same data/split or evaluated with
released weights, and the number of context pairs $K$ used at
inference}.

\noindent\textbf{Metrics.} Segmentation accuracy is reported as Dice and
normalized surface distance (NSD); compute is reported as wall-clock
inference time, FLOPs, and peak VRAM.

\subsection{Fixed-Spacing Comparison}
\label{sec:fixed-spacing}

We first compare all methods at a fixed 3\,mm isotropic spacing on
TotalSegmentator CT crops, isolating architectural differences in
image--label interaction (\cref{sec:biaxial}) from the coarse-to-fine
cascade evaluated separately in \cref{sec:c2f-ablation}.

\begin{table}[t]
\centering
\caption{Dice, NSD, and compute at fixed 3\,mm spacing on
TotalSegmentator CT held-out classes. \TODO{fill in.}}
\label{tab:fixed-spacing}
\small
\begin{tabular}{@{}lcccc@{}}
\toprule
Method & Dice & NSD & Time (s) & VRAM (GB) \\
\midrule
Iris & \TODO{} & \TODO{} & \TODO{} & \TODO{} \\
Medverse & \TODO{} & \TODO{} & \TODO{} & \TODO{} \\
PatchICL-3D (ours) & \TODO{} & \TODO{} & \TODO{} & \TODO{} \\
\bottomrule
\end{tabular}
\end{table}

\TODO{Discuss overall numbers, then break down by anatomical category
(organs, bones, vessels, muscles) to identify where bi-axial attention
helps most -- we expect an advantage on small/thin structures, similar
to the pattern observed for patch-based selection in our 2D work.}

\subsection{Coarse-to-Fine Ablations}
\label{sec:c2f-ablation}

We ablate the three cascade mechanisms of \cref{sec:cascade,sec:registers}
independently: (i) the query image prior
(\cref{eq:query-prior}) vs.\ a blank label-token initialization at every
level; (ii) region-restricted fine processing (sliding window over the
previous level's predicted zone) vs.\ a full-volume sliding window at
every level; and (iii) register-carried cross-level training
(\cref{eq:register-carry}) vs.\ independently-supervised levels
connected only through the rendered query prior.

\label{sec:cross-level-ablation}
\begin{table}[t]
\centering
\caption{Effect of each cascade component on Dice and compute.
\TODO{fill in; one row per ablation, cumulative or leave-one-out.}}
\label{tab:c2f-ablation}
\small
\begin{tabular}{@{}lccc@{}}
\toprule
Configuration & Dice & Time (s) & FLOPs \\
\midrule
Full model & \TODO{} & \TODO{} & \TODO{} \\
$-$ query prior & \TODO{} & \TODO{} & \TODO{} \\
$-$ region restriction & \TODO{} & \TODO{} & \TODO{} \\
$-$ register carry & \TODO{} & \TODO{} & \TODO{} \\
\bottomrule
\end{tabular}
\end{table}

\TODO{Report the accuracy/compute Pareto trade-off as sliding-window
restriction is relaxed (full volume $\rightarrow$ tight crop around the
previous prediction), analogous to the resolution/FLOPs Pareto plot in
our 2D paper~\cite{camaretPatchICL2026}.}

\subsection{Generalization}
\label{sec:generalization}

To test how well PatchICL-3D's accuracy/compute trade-off transfers
beyond the training distribution, we evaluate on three settings of
increasing distribution shift.

\noindent\textbf{Other CT datasets.} \TODO{list held-out CT cohorts and
results.}

\noindent\textbf{Other modalities.} \TODO{MRI (e.g.\ TotalSegmentator
MRI~\cite{dantonoliTotalSegmentatorMRIRobust2025}) evaluated with models
trained only on CT.}

\noindent\textbf{Far out-of-distribution tasks.} \TODO{non-anatomical or
pathology-driven segmentation tasks probing whether the region-restricted
cascade still localizes correctly when the coarse level's prior is
unreliable.}
